%% =====================================================================
%%  4 Physical Model
%% =====================================================================
%%  GUIDE
%%   - Build in layers: layer 1 the minimum model, then add one effect per
%%     subsection, saying what changed in the answer (many gold papers use
%%     explicit "4.1 free flow / 4.2 atmosphere / 4.3 drag" layering).
%%   - State the governing equation for each layer, its regime of validity,
%%     and number it if you refer to it later.
%%   - Derive or motivate; do not drop equations from nowhere.
%%   - After each equation, define every symbol and check the dimensions
%%     once in words, e.g. "(Pa/s on both sides)".
%% =====================================================================
\section{Physical Model}
\label{sec:model}

\subsection{Layer 1: minimum model, isentropic discharge}
\label{sec:model-basic}

The module is a reservoir of volume $V$ at uniform pressure $p(t)$; the
breach is an orifice of area $A=\pi d^2/4$. Conservation of mass for the
reservoir, with the outflow evaluated at the orifice, gives
\begin{equation}
  \frac{\dd p}{\dd t} = -C_d\,\frac{\gamma p A}{V}\,u(p,T),
  \label{eq:mass}
\end{equation}
where $u$ is the exit speed, $\gamma$ the heat-capacity ratio and $C_d$ a
discharge coefficient. Both sides of \eqref{eq:mass} carry \si{Pa/s}. The
coefficient absorbs the vena-contracta contraction and the ratio of
orifice to reservoir density; it is the single free parameter of the model
and is calibrated in Section~\ref{sec:disc-validation}, not taken from a
handbook.

For an ideal gas expanding from rest, the exit speed is
\begin{equation}
  u^2 = \frac{2\gamma}{\gamma-1}RT
  \left[1-\left(\frac{p_a}{p}\right)^{(\gamma-1)/\gamma}\right],
  \label{eq:subsonic}
\end{equation}
which holds until the flow chokes. Choking occurs when the bracket reaches
its maximum, at $p/p_a=((\gamma+1)/2)^{\gamma/(\gamma-1)}=1.89$ for air;
beyond that the exit Mach number stays at unity and the speed is set by the
reservoir state alone,
\begin{equation}
  u_{\rm ch} = \sqrt{\frac{2\gamma}{\gamma+1}RT}.
  \label{eq:choked}
\end{equation}
For a vacuum exterior, $p_a\to0$ and the ratio $p/p_a$ never falls below
1.89, so the flow is choked from the first instant to the last and
\eqref{eq:subsonic} is never selected. We keep both branches in the code
because a real breach may vent into a partially pressurised compartment,
but for the nominal case only \eqref{eq:choked} is active. Equations
\eqref{eq:mass} and \eqref{eq:choked} are the minimum model: one ordinary
differential equation for $p(t)$, with $T$ held at $T_0$.

\subsection{Layer 2: wall heat transfer}
\label{sec:model-thermal}

Adiabatic expansion cools the gas, and the cooling is visible in the
solution: at $d=\SI{10}{cm}$ the model above ends at \SI{208}{K}. The
walls, however, hold the module near its initial temperature on a time
scale of a few hundred seconds, so the gas is warmed as it expands. A
lumped first-order relaxation captures this,
\begin{equation}
  \frac{\dd T}{\dd t}
  = \frac{\gamma-1}{\gamma}\frac{T}{p}\frac{\dd p}{\dd t}
  - \frac{1}{\tau_w}\left(T-T_w\right),
  \label{eq:temp}
\end{equation}
where the first term is the isentropic relation $T\propto
p^{(\gamma-1)/\gamma}$ written as a rate, and $\tau_w$ is the wall
relaxation time estimated from a lumped-capacitance balance. For the
nominal module $\tau_w\approx\SI{300}{s}$, comparable with $t_{0.3}$, so
the correction is not negligible: it lowers $t_{0.3}$ by \SI{2}{\percent}
at $d=\SI{10}{cm}$ and by \SI{5}{\percent} at \SI{5}{cm}, because a warmer
gas is less dense and carries less mass out of the same hole.

\subsection{Layer 3: dimensionless form}
\label{sec:model-pi}

With $P=p/p_0$ and $\theta=t/\tau$, where
\begin{equation}
  \tau = \frac{V}{\gamma C_d A c_s},\qquad c_s=\sqrt{\gamma R T_0},
  \label{eq:tau}
\end{equation}
the system \eqref{eq:mass}--\eqref{eq:temp} depends on two groups only,
\begin{equation}
  \frac{\dd P}{\dd\theta} = -P\,\frac{u(P)}{c_s},\qquad
  \Pi_T = \frac{\tau}{\tau_w},
  \label{eq:pi}
\end{equation}
so one numerical solution covers every module size and every hole
diameter, provided $\gamma$ and $\Pi_T$ are unchanged. For the nominal
case $\tau=\SI{203}{s}$, and $t_{0.3}/\tau=1.41$. Because $A\propto d^2$,
\eqref{eq:tau} already implies $t_{0.3}\propto d^{-2}$; the measured
exponent of Section~\ref{sec:results} is slightly shallower because
$\Pi_T$ changes with $d$.
